\documentclass[12pt,a4paper]{article}

\usepackage[a4paper,margin=2.2cm]{geometry}
\usepackage{iftex}

\ifPDFTeX
  \usepackage[utf8]{inputenc}
  \usepackage[T5]{fontenc}
  \usepackage[vietnamese]{babel}
\else
  \usepackage{fontspec}
  \usepackage{polyglossia}
  \setdefaultlanguage{vietnamese}
  \setmainfont{TeX Gyre Termes}
\fi

\usepackage{setspace}
\usepackage{graphicx}
\usepackage{longtable}
\usepackage{booktabs}
\usepackage{array}
\usepackage{enumitem}
\usepackage{xcolor}
\usepackage[hidelinks]{hyperref}

\setstretch{1.25}
\setlist[itemize]{leftmargin=*}
\setlist[enumerate]{leftmargin=*}

\title{\textbf{BÁO CÁO BÀI TẬP LỚN LẬP TRÌNH WEB}\\
HK2 2025--2026\\
\large Thiết kế giao diện và xây dựng các tính năng cơ bản cho website công ty -- doanh nghiệp}
\author{Nhóm: Team36\\
Email liên hệ: team36-web@hcmut.edu.vn}
\date{Tháng 04 năm 2026}

\begin{document}

\maketitle
\tableofcontents
\newpage

\section{Giới thiệu đề tài}

\subsection{Bối cảnh}
Đề tài yêu cầu xây dựng website công ty -- doanh nghiệp bằng các công nghệ đã học: HTML5, CSS3, Javascript, PHP và PHP \& MySQL. Hệ thống phải đáp ứng các vai trò người dùng gồm khách, thành viên, quản trị viên và tuân thủ các ràng buộc của môn học (không dùng PHP framework/CMF/CMS).

\subsection{Mục tiêu thực hiện}
\begin{itemize}
  \item Chuẩn hóa các trang giao diện public theo đúng đề.
  \item Hoàn thiện luồng người dùng: khách, thành viên, quản trị viên.
  \item Đảm bảo các tính năng bắt buộc: tìm kiếm, phân trang, quản lý liên hệ, quản lý người dùng.
  \item Trình bày báo cáo bám sát mã nguồn thực tế, không nêu tính năng chưa có code.
\end{itemize}

\section{Công nghệ và kiến trúc}

\subsection{Công nghệ sử dụng}
\begin{itemize}
  \item Frontend: HTML5, CSS3, Bootstrap 4, Javascript.
  \item Backend: PHP thuần theo kiểu MVC nhẹ.
  \item Database: SQLite cho môi trường demo; có tệp tham chiếu cấu trúc MySQL trong database.sql.
\end{itemize}

\subsection{Mô hình tổ chức mã nguồn}
\begin{itemize}
  \item Model: thao tác dữ liệu (User, Product, News, Contact).
  \item View: giao diện public, tài khoản, quản trị.
  \item Controller: điều phối request qua tệp index.php.
\end{itemize}

\section{Thiết kế cơ sở dữ liệu tóm tắt}
\begin{longtable}{p{0.24\textwidth}p{0.71\textwidth}}
\toprule
Bảng & Vai trò \\
\midrule
users & Lưu tài khoản, vai trò, trạng thái, avatar \\
products & Lưu thông tin sản phẩm/dịch vụ, giá, mô tả, hình ảnh \\
news & Lưu bài viết tin tức \\
contacts & Lưu liên hệ khách hàng và trạng thái xử lý (new/read/replied) \\
\bottomrule
\end{longtable}

\section{Các tính năng đã hiện thực}

\subsection{Nhóm khách}
\begin{itemize}
  \item Xem các trang public: trang chủ, giới thiệu, sản phẩm, bảng giá, liên hệ, FAQ, tin tức.
  \item Tìm kiếm sản phẩm theo từ khóa.
  \item Tìm kiếm tin tức theo từ khóa.
  \item Xem trang chi tiết sản phẩm và chi tiết tin tức.
  \item Gửi form liên hệ (có kiểm tra dữ liệu đầu vào).
\end{itemize}

\subsection{Nhóm thành viên}
\begin{itemize}
  \item Đăng ký, đăng nhập.
  \item Cập nhật thông tin cá nhân.
  \item Cập nhật avatar.
  \item Đổi mật khẩu (có kiểm tra mật khẩu cũ và xác nhận mật khẩu mới).
\end{itemize}

\subsection{Nhóm quản trị viên}
\begin{itemize}
  \item Quản lý sản phẩm: xem, thêm, xóa.
  \item Quản lý thành viên: xem danh sách, khóa/mở khóa, reset mật khẩu, xóa.
  \item Quản lý liên hệ khách hàng: cập nhật trạng thái new/read/replied.
  \item Phân trang danh sách liên hệ trong khu vực admin.
\end{itemize}

\subsection{Phân trang}
\begin{itemize}
  \item Danh sách sản phẩm.
  \item Danh sách tin tức.
  \item Danh sách liên hệ trong admin.
\end{itemize}

\section{Bảo mật cơ bản và SEO cơ bản}

\subsection{Bảo mật cơ bản}
\begin{itemize}
  \item Mật khẩu được băm bằng password\_hash và xác thực bằng password\_verify.
  \item Escape output bằng htmlspecialchars để giảm rủi ro XSS.
  \item Dùng prepared statement khi truy vấn cơ sở dữ liệu.
  \item Kiểm soát quyền theo role ở khu vực quản trị.
\end{itemize}

\subsection{SEO cơ bản}
\begin{itemize}
  \item Thiết lập title theo từng trang.
  \item Cấu trúc nội dung có heading rõ ràng.
  \item URL có tham số phân biệt chức năng tìm kiếm và phân trang.
\end{itemize}

\section{Môi trường chạy và cách cài đặt}

\subsection{Yêu cầu môi trường}
\begin{itemize}
  \item PHP từ phiên bản 7.0 trở lên (khuyến nghị PHP 8.x).
  \item Web server Apache/Nginx hoặc chạy nhanh bằng PHP built-in server.
  \item Cơ sở dữ liệu SQLite hoặc MySQL theo kịch bản triển khai.
\end{itemize}

\subsection{Cách chạy nhanh}
\begin{enumerate}
  \item Vào thư mục Assignment1 - Team36.
  \item Chạy lệnh php -S localhost:8000.
  \item Mở trình duyệt tại đường dẫn /index.php.
  \item Tài khoản admin mẫu: admin@vng.com.
\end{enumerate}

\section{Đối chiếu theo yêu cầu đề bài}

\begin{longtable}{p{0.40\textwidth}p{0.20\textwidth}p{0.35\textwidth}}
\toprule
Yêu cầu đề bài & Trạng thái & Ghi chú theo mã nguồn hiện tại \\
\midrule
Trang public cơ bản (home/about/products/pricing/contact/faq/news) & Đã có & Có routing và giao diện tương ứng cho từng trang. \\
Đăng ký/Đăng nhập bằng PHP + CSDL & Đã có & Luồng hoạt động ổn định, mật khẩu băm đúng chuẩn. \\
Cập nhật hồ sơ, mật khẩu, avatar & Đã có & Có luồng cập nhật hồ sơ và upload avatar. \\
Tìm kiếm sản phẩm/tin tức & Đã có & Tìm kiếm theo từ khóa ở phần public. \\
Phân trang dữ liệu dài & Đã có một phần & Đã có cho products, news, contacts; chưa đồng bộ toàn bộ trang quản trị. \\
Quản lý user cho admin (xem/sửa/trạng thái/reset/xóa) & Đã có một phần & Đã có xem + khóa/mở + reset + xóa; chưa có chỉnh sửa chi tiết hồ sơ user trong admin. \\
Quản lý liên hệ khách hàng & Đã có & Có xem danh sách và cập nhật trạng thái xử lý. \\
Quản lý sản phẩm (xem/thêm/sửa/xóa) & Đã có một phần & Đã có xem + thêm + xóa; còn thiếu sửa đầy đủ. \\
Quản lý tin tức admin (xem/thêm/sửa/xóa + metadata) & Còn thiếu & Chưa có module quản trị tin tức đầy đủ. \\
Bình luận/đánh giá của thành viên & Còn thiếu & Chưa có bảng dữ liệu và giao diện tương ứng. \\
Giỏ hàng/đơn hàng (nếu áp dụng) & Còn thiếu & Chưa triển khai trong phiên bản hiện tại. \\
\bottomrule
\end{longtable}

\section{Hạng mục ưu tiên bổ sung}

\begin{enumerate}
  \item Hoàn thiện CRUD tin tức phía admin (xem, thêm, sửa, xóa, tìm kiếm, phân trang).
  \item Bổ sung module bình luận/đánh giá có duyệt nội dung cơ bản.
  \item Nếu theo hướng thương mại điện tử: bổ sung giỏ hàng, đơn hàng và trạng thái đơn.
  \item Chuẩn hóa dashboard quản trị theo template Srtdash đúng tinh thần đề.
\end{enumerate}

\section{Phân công công việc nhóm}
\begin{itemize}
  \item Thành viên 1: giao diện trang chủ, liên hệ; luồng xử lý liên hệ.
  \item Thành viên 2: giao diện giới thiệu, hỏi đáp; responsive.
  \item Thành viên 3: danh sách sản phẩm, tìm kiếm, phân trang, chi tiết sản phẩm.
  \item Thành viên 4: danh sách tin tức, tìm kiếm, phân trang, chi tiết tin tức.
  \item Phần chung: đăng ký/đăng nhập, phân quyền, hồ sơ người dùng, module quản trị cơ bản.
\end{itemize}

\section{Kết luận}
Hệ thống đã hoàn thành phần nền tảng quan trọng theo yêu cầu môn học: giao diện public, xác thực người dùng, hồ sơ cá nhân, tìm kiếm, phân trang cơ bản và quản trị liên hệ/người dùng ở mức cơ bản. Để đạt mức điểm cao hơn theo đề, nhóm cần tập trung hoàn thiện các phần còn thiếu gồm quản trị tin tức đầy đủ, bình luận/đánh giá và giỏ hàng/đơn hàng (nếu áp dụng).

\section*{Tài liệu tham khảo}
\begin{itemize}
  \item Đề bài Lập trình Web HK2 2025--2026 (BKel).
  \item PHP Documentation: \url{https://www.php.net/docs.php}
  \item SQLite Documentation: \url{https://www.sqlite.org/docs.html}
  \item Bootstrap 4 Documentation: \url{https://getbootstrap.com/docs/4.0/getting-started/introduction/}
\end{itemize}

\end{document}
